\documentclass[a4paper]{article}

\usepackage{amsmath}
\usepackage{amssymb}
\usepackage{array}
\usepackage{graphicx}
\usepackage{hyperref}
\usepackage{listings}
\usepackage{longtable}
\usepackage{natbib}
\usepackage{nicefrac}
\usepackage{ulem}

\newcommand{\pequod}{PEQUOD}
\newcommand{\pequodexpanded}{Parallel Quasi-Geostrophic Model}
\newcommand{\version}{0.9.5-dev}

\begin{document}

\title{Building \pequod\ \version\ dependencies}
\maketitle

\section{Introduction}

This document gives step-by-step instructions for a ground-up build of \pequod,
including a build of all model dependencies. It is assumed that a minimal build
environment is available, including versions of Bash, the GCC C and C++ compilers, m4,
and GNU make, as well as a number of standard command line utilities such as
tar, Wget and md5sum. These steps have been tested on a minimal vanilla
installation of the 64-bit version of Ubuntu 12.04 (Precise
Pangolin)\footnote{Available from \url{http://www.ubuntu.com}.}, with a small
number of additional packages installed via:
\begin{lstlisting}[language = bash, frame = single]
sudo apt-get install g++ gcc m4 make
\end{lstlisting}

The instructions in this document download and compile free and open source
versions of required \pequod\ dependencies. Clusters generally include
recommended compilers and optimised libraries, and hence these instructions are
not intended for use on HPC resources.

The steps in this document should be followed in order, and a successful build
of each dependency should be achieved before moving on to the next section. All
instructions in this document assume usage of the Bash shell. If using an
alternative shell, switch to Bash using the command:
\begin{lstlisting}[language = sh, frame = single]
bash
\end{lstlisting}

\section{Working directory}\label{sect:dir}

First, create a working directory:
\begin{lstlisting}[language = bash, frame = single]
mkdir -pv $HOME/build/pequod-dependencies
export PDEP_DIR=$HOME/build/pequod-dependencies
\end{lstlisting}
The first command creates the directory \verb+$HOME/build/pequod-dependencies+,
although you can replace this with any target directory. The second sets an
environment variable \verb+PDEP_DIR+ pointing at this directory. Note that you
will need to reset the environment variable \verb+PDEP_DIR+ if you terminate the
shell session. Next create a download directory, a build directory, and an
installation directory with:
\begin{lstlisting}[language = bash, frame = single]
mkdir -v $PDEP_DIR/src
mkdir -v $PDEP_DIR/build
mkdir -v $PDEP_DIR/install
\end{lstlisting}

Now set the following environment variables:
\begin{lstlisting}[language = bash, frame = single]
export LIBRARY_PATH=$PDEP_DIR/install/lib:\
$PDEP_DIR/install/lib64
export LD_LIBRARY_PATH=$PDEP_DIR/install/lib:\
$PDEP_DIR/install/lib64
\end{lstlisting}
Note that you will need to reset these environment variables if you
terminate the shell session.

\section{libgmp}

The GNU Multiple Precision Arithmetic library is available from
\url{http://gmplib.org/}. Download it using:
\begin{lstlisting}[language = bash, frame = single]
cd $PDEP_DIR/src
wget ftp://ftp.gmplib.org/pub/gmp-5.0.5/gmp-5.0.5.tar.bz2
\end{lstlisting}
Verify the download with:
\begin{lstlisting}[language = bash, frame = single]
md5sum gmp-5.0.5.tar.bz2
\end{lstlisting}
The correct checksum is \verb+041487d25e9c230b0c42b106361055fe+. Extract the
source files using:
\begin{lstlisting}[language = bash, frame = single]
tar -xvf gmp-5.0.5.tar.bz2
\end{lstlisting}
Create and change to a build directory:
\begin{lstlisting}[language = bash, frame = single]
mkdir -v $PDEP_DIR/build/gmp
cd $PDEP_DIR/build/gmp
\end{lstlisting}
Now configure, compile, and install the library with:
\begin{lstlisting}[language = bash, frame = single]
$PDEP_DIR/src/gmp-5.0.5/configure \
  --prefix=$PDEP_DIR/install
make
make install
\end{lstlisting}

\section{libmpfr}

The GNU MPFR library is available from \url{http://www.mpfr.org/}. Download it
using:
\begin{lstlisting}[language = bash, frame = single]
cd $PDEP_DIR/src
wget http://www.mpfr.org/mpfr-current/mpfr-3.1.1.tar.bz2
\end{lstlisting}
Verify the download with:
\begin{lstlisting}[language = bash, frame = single]
md5sum mpfr-3.1.1.tar.bz2
\end{lstlisting}
The correct checksum is \verb+e90e0075bb1b5f626c6e31aaa9c64e3b+. Extract the
source files using:
\begin{lstlisting}[language = bash, frame = single]
tar -xvf mpfr-3.1.1.tar.bz2
\end{lstlisting}
Create and change to a build directory:
\begin{lstlisting}[language = bash, frame = single]
mkdir -v $PDEP_DIR/build/mpfr
cd $PDEP_DIR/build/mpfr
\end{lstlisting}
Now configure, compile, and install the library with:
\begin{lstlisting}[language = bash, frame = single]
$PDEP_DIR/src/mpfr-3.1.1/configure \
  --prefix=$PDEP_DIR/install --with-gmp=$PDEP_DIR/install
make
make install
\end{lstlisting}

\section{libmpc}

Mpc is available from \url{http://www.multiprecision.org/}. Download it using:
\begin{lstlisting}[language = bash, frame = single]
cd $PDEP_DIR/src
wget http://www.multiprecision.org/mpc/download/\
mpc-1.0.1.tar.gz
\end{lstlisting}
Verify the download with:
\begin{lstlisting}[language = bash, frame = single]
md5sum mpc-1.0.1.tar.gz
\end{lstlisting}
The correct checksum is \verb+b32a2e1a3daa392372fbd586d1ed3679+. Extract the
source files using:
\begin{lstlisting}[language = bash, frame = single]
tar -xvf mpc-1.0.1.tar.gz
\end{lstlisting}
Create and change to a build directory:
\begin{lstlisting}[language = bash, frame = single]
mkdir -v $PDEP_DIR/build/mpc
cd $PDEP_DIR/build/mpc
\end{lstlisting}
Now configure, compile, and install the library with:
\begin{lstlisting}[language = bash, frame = single]
$PDEP_DIR/src/mpc-1.0.1/configure \
  --prefix=$PDEP_DIR/install \
  --with-gmp=$PDEP_DIR/install \
  --with-mpfr=$PDEP_DIR/install
make
make install
\end{lstlisting}

\section{GCC}

The GNU Compiler Collection is available from \url{http://gcc.gnu.org/}.
Download it using:
\begin{lstlisting}[language = bash, frame = single]
cd $PDEP_DIR/src
wget http://mirrors.kernel.org/gnu/gcc/gcc-4.7.1/\
gcc-4.7.1.tar.bz2
\end{lstlisting}
Verify the download with:
\begin{lstlisting}[language = bash, frame = single]
md5sum gcc-4.7.1.tar.bz2
\end{lstlisting}
The correct checksum is \verb+933e6f15f51c031060af64a9e14149ff+. Extract the
source files using:
\begin{lstlisting}[language = bash, frame = single]
tar -xvf gcc-4.7.1.tar.bz2
\end{lstlisting}
Create and change to a build directory:
\begin{lstlisting}[language = bash, frame = single]
mkdir -v $PDEP_DIR/build/gcc
cd $PDEP_DIR/build/gcc
\end{lstlisting}
Now configure, compile, and install GCC with:
\begin{lstlisting}[language = bash, frame = single]
$PDEP_DIR/src/gcc-4.7.1/configure \
  --prefix=$PDEP_DIR/install \
  --with-gmp=$PDEP_DIR/install \
  --with-mpfr=$PDEP_DIR/install \
  --with-mpc=$PDEP_DIR/install \
  --enable-languages=c,c++,fortran --enable-bootstrap \
  --disable-multilib
make boostrap
make install
\end{lstlisting}
Note that the \verb+make bootstrap+ command may take a long time (of order
hours) to complete.

\subsection{Known issues}

\subsubsection*{64-bit Ubuntu 12.04 (Precise Pangolin)}

The compilation fails to find a header file \verb+bits/predefs.h+. This can be
worked around by defining environment variables:
\begin{lstlisting}[language = bash, frame = single]
export C_INCLUDE_PATH=/usr/include/x86_64-linux-gnu
export CPLUS_INCLUDE_PATH=/usr/include/x86_64-linux-gnu
\end{lstlisting}
The compilation further fails to find a file \verb+crti.o+. This can be worked
around by updating the \verb+LIBRARY_PATH+ environment variable to:
\begin{lstlisting}[language = bash, frame = single]
export LIBRARY_PATH=$PDEP_DIR/install/lib:\
$PDEP_DIR/install/lib64:/usr/lib/x86_64-linux-gnu
\end{lstlisting}
The environment variables should be defined before the \verb+make bootstrap+
command is executed. The environment variables should be thus defined when
building all following dependencies.

\section{Open MPI}

Open MPI is available from \url{http://www.open-mpi.org/}. Download it using:
\begin{lstlisting}[language = bash, frame = single]
cd $PDEP_DIR/src
wget http://www.open-mpi.org/software/ompi/v1.6/\
downloads/openmpi-1.6.2.tar.bz2
\end{lstlisting}
Verify the download with:
\begin{lstlisting}[language = bash, frame = single]
md5sum openmpi-1.6.2.tar.bz2
\end{lstlisting}
The correct checksum is \verb+219754715a8e40beb468bbc8f0b3251a+. Extract the
source files using:
\begin{lstlisting}[language = bash, frame = single]
tar -xvf openmpi-1.6.2.tar.bz2
\end{lstlisting}
Create and change to a build directory:
\begin{lstlisting}[language = bash, frame = single]
mkdir -v $PDEP_DIR/build/openmpi
cd $PDEP_DIR/build/openmpi
\end{lstlisting}
Now configure, compile, and install Open MPI with:
\begin{lstlisting}[language = bash, frame = single]
CC=$PDEP_DIR/install/bin/gcc \
  CPP=$PDEP_DIR/install/bin/cpp \
  CXX=$PDEP_DIR/install/bin/g++ \
  CXXPP=$PDEP_DIR/install/bin/cpp \
  F77=$PDEP_DIR/install/bin/gfortran \
  FC=$PDEP_DIR/install/bin/gfortran \
  $PDEP_DIR/src/openmpi-1.6.2/configure \
  --prefix=$PDEP_DIR/install
make
make install
\end{lstlisting}

\section{liblapack}

An implementation of LAPACK is available from \url{http://www.netlib.org/lapack/}.
Download it using:
\begin{lstlisting}[language = bash, frame = single]
cd $PDEP_DIR/src
wget http://www.netlib.org/lapack/lapack-3.4.1.tgz
\end{lstlisting}
Verify the download with:
\begin{lstlisting}[language = bash, frame = single]
md5sum lapack-3.4.1.tgz
\end{lstlisting}
The correct checksum is \verb+44c3869c38c8335c2b9c2a8bb276eb55+. Extract the
source files using:
\begin{lstlisting}[language = bash, frame = single]
tar -xvf lapack-3.4.1.tgz
\end{lstlisting}
Create and change to a build directory:
\begin{lstlisting}[language = bash, frame = single]
cp -rv $PDEP_DIR/src/lapack-3.4.1 $PDEP_DIR/build/lapack
cd $PDEP_DIR/build/lapack
\end{lstlisting}
Create a \verb+make.inc+ file:
\begin{lstlisting}[language = bash, frame = single]
cp -v make.inc.example make.inc
\end{lstlisting}
Edit the \verb+make.inc+ file so that the following variables have the indicated
values:
\begin{lstlisting}[language = bash, frame = single]
FORTRAN  = $(PDEP_DIR)/install/bin/mpif90
LOADER   = $(PDEP_DIR)/install/bin/mpif90
CC = $(PDEP_DIR)/install/bin/mpicc
\end{lstlisting}
Compile the library with:
\begin{lstlisting}[language = bash, frame = single]
make blaslib all
\end{lstlisting}
Install the library with:
\begin{lstlisting}[language = bash, frame = single]
cp -v librefblas.a $PDEP_DIR/install/lib/libblas.a
cp -v liblapack.a $PDEP_DIR/install/lib
\end{lstlisting}

\section{Compiling \pequod}

In order to compile \pequod\ the following environment variables must be set:
\begin{lstlisting}[language = bash, frame = single]
export MPI_DIR=$PDEP_DIR/install
export LD_LIBRARY_PATH=$PDEP_DIR/install/lib:\
$PDEP_DIR/install/lib64
export LIBRARY_PATH=$PDEP_DIR/install/lib:\
$PDEP_DIR/install/lib64
\end{lstlisting}

These can be set in your \verb+.bashrc+ file by adding the lines:
\begin{lstlisting}[language = bash, frame = single]
export MPI_DIR=$HOME/build/pequod-dependencies/install
if [ -z "${LD_LIBRARY_PATH}" ]; then
  export LD_LIBRARY_PATH=\
$HOME/build/pequod-dependencies/install/lib:\
$HOME/build/pequod-dependencies/install/lib64
else
  export LD_LIBRARY_PATH=\
$HOME/build/pequod-dependencies/install/lib:\
$HOME/build/pequod-dependencies/install/lib64:\
$LD_LIBRARY_PATH
fi
if [ -z "${LIBRARY_PATH}" ]; then
  export LIBRARY_PATH=\
$HOME/build/pequod-dependencies/install/lib:\
$HOME/build/pequod-dependencies/install/lib64
else
  export LIBRARY_PATH=\
$HOME/build/pequod-dependencies/install/lib:\
$HOME/build/pequod-dependencies/install/lib64:\
$LIBRARY_PATH
fi
\end{lstlisting}
The directory \verb+$HOME/build/pequod-dependencies+ should be replaced with the
directory created in section \ref{sect:dir}.

In the \pequod\ source directory, edit \verb+makefile.gfortran+ and remove the
\verb+-floop-interchange+, \verb+-floop-strip-mine+, and \verb+-floop-block+
flags from the \verb+OPTIMISATION+ variable. The \verb+LIBS+ variable should be
set to:
\begin{lstlisting}[language = bash, frame = single]
LIBS = -llapack -lblas
\end{lstlisting}
Ensure that \verb+makefile.compiler+ includes \verb+makefile.gfortran+.

\pequod\ can now be configured and compiled as described in the getting
started guide and the \pequod\ manual. Note that, when running in parallel,
if \verb+$PDEP_DIR/install/bin+ is not included in the \verb+PATH+ environment
variable then you should include the full path to \verb+mpi_run+. For example,
when running \pequod\ on 4 processes, use the command:
\begin{lstlisting}[language = bash, frame = single]
$MPI_DIR/bin/mpirun -np 4 ./qg3
\end{lstlisting}

\subsection{Known issues}

\subsubsection*{64-bit Ubuntu 12.04 (Precise Pangolin)}

The compilation fails to find a file \verb+crti.o+. This can be worked
around by setting the \verb+LIBRARY_PATH+ environment variable to:
\begin{lstlisting}[language = bash, frame = single]
export LIBRARY_PATH=$PDEP_DIR/install/lib:\
$PDEP_DIR/install/lib64:/usr/lib/x86_64-linux-gnu
\end{lstlisting}
If modifying your \verb+.bashrc+ file, add:
\begin{lstlisting}[language = bash, frame = single]
export LIBRARY_PATH=$LIBRARY_PATH:\
/usr/lib/x86_64-linux-gnu
\end{lstlisting}
after all the lines indicated above.

%\bibliography{bibliography}
%\bibliographystyle{plainnat}

\end{document}
